\documentclass[11pt]{article}
\usepackage[T1]{fontenc}
\usepackage{lmodern}
\usepackage{amsmath,amssymb,amsthm}
\usepackage{graphicx}
\usepackage[margin=1in]{geometry}
\usepackage{microtype}
\usepackage[colorlinks=true,linkcolor=blue,citecolor=blue,urlcolor=blue]{hyperref}

\newtheorem{theorem}{Theorem}
\newtheorem{lemma}[theorem]{Lemma}
\newtheorem{corollary}[theorem]{Corollary}
\newcommand{\Lcal}{\mathcal L}
\newcommand{\Dom}{\operatorname{Dom}}
\newcommand{\R}{\mathbb R}

\title{The subcritical M\'etivier Schr\"odinger potential is essentially self-adjoint}
\author{Candidate solution packet for arXiv:1610.09850}
\date{13 August 2026}

\begin{document}
\maketitle

\begin{center}
\fbox{\parbox{0.91\linewidth}{\textbf{Status.} Candidate full affirmative
answer to Remark 4.8, likely valid, awaiting specialist review.}}
\end{center}

\section{Source question}

Bruno and Calzi ask in Remark 4.8 of \cite{BC}, on PDF page 12, whether the
Schr\"odinger operator associated with their weighted sub-Laplacian is
essentially self-adjoint in the range $0<\alpha<2$.

\begin{center}
\includegraphics[width=\linewidth]{figures/open_problem_crop.png}
\end{center}

We answer yes throughout the open range.  The negative part of the potential
is compactly supported and strictly milder than the homogeneous inverse-square
scale.  A quantitative resolvent estimate therefore turns the Kato inequality
already proved in the source into a strict contraction of the deficiency
space.

\section{A general criterion}

Let $G$ be a stratified group of homogeneous dimension $Q\geq4$, let
$\Lcal=-\sum_jX_j^2$ be its positive sub-Laplacian, and let Haar measure be
fixed.  We use the self-adjoint closure of $(\Lcal,C_c^\infty(G))$.

\begin{theorem}\label{thm:criterion}
Let $V\colon G\to\R$ belong to $L^2_{\mathrm{loc}}(G)$.  If its negative
part $W=V_-:=\max\{-V,0\}$ belongs to $L^p(G)$ for some $p>Q/2$, then
\[
 (\Lcal+V,C_c^\infty(G))
\]
is essentially self-adjoint on $L^2(G)$.
\end{theorem}

The global $L^p$ hypothesis is convenient rather than optimal.  In the
application below $W$ is compactly supported.

\section{Proof intuition}

The source proves essential self-adjointness when $V$ is bounded below by
applying Kato's inequality to a deficiency vector.  Here $V$ tends to
$-\infty$ at the identity, so that argument cannot discard the term $V_-|f|$.
Instead, apply the free resolvent $R_\lambda=(\Lcal+\lambda)^{-1}$.  Heat-kernel
scaling gives
\[
 \|W R_\lambda\|_{2\to2}
 \leq C\|W\|_p\lambda^{-1+Q/(2p)}=o(1).
\]
Because $p>Q/2$, this norm is strictly less than one for large $\lambda$.
Kato's inequality says that a deficiency vector satisfies
$(\Lcal+\lambda)|f|\leq W|f|$.  Testing against $R_\lambda\psi$ then yields
$\|f\|_2\leq\|W R_\lambda\|_{2\to2}\|f\|_2$, forcing $f=0$.

\section{Resolvent estimate}

\begin{lemma}\label{lem:resolvent}
If $p>Q/2$ and $W\in L^p(G)$, then, for $\lambda>0$,
\begin{equation}\label{eq:resolvent}
 \|W(\Lcal+\lambda)^{-1}\|_{2\to2}
 \leq C_{G,p}\|W\|_p\lambda^{-1+Q/(2p)}.
\end{equation}
\end{lemma}

\begin{proof}
Put $r=2p/(p-2)$, so that $1/2=1/p+1/r$.  The sub-Laplacian heat kernel has
the homogeneous scaling
\[
 h_t(x)=t^{-Q/2}h_1(\delta_{t^{-1/2}}x).
\]
Young's convolution inequality consequently gives
\[
 \|e^{-t\Lcal}\|_{2\to r}
 \leq C t^{-\frac Q2(\frac12-\frac1r)}
 =C t^{-Q/(2p)}.
\]
Using the semigroup representation of the resolvent,
\[
 \|(\Lcal+\lambda)^{-1}\|_{2\to r}
 \leq C\int_0^\infty e^{-\lambda t}t^{-Q/(2p)}\,dt
 =C_{G,p}\lambda^{-1+Q/(2p)}.
\]
H\"older's inequality for multiplication by $W$ proves
\eqref{eq:resolvent}.  The integral is finite precisely because $p>Q/2$.
\end{proof}

\section{Proof of the criterion}

First, the first-order homogeneous Sobolev inequality and interpolation imply,
for $s=2p/(p-1)<2Q/(Q-2)$ and every $\eta>0$,
\[
 \|u\|_s^2
 \leq \eta\|\Lcal^{1/2}u\|_2^2+C_\eta\|u\|_2^2,
 \qquad u\in C_c^\infty(G).
\]
Thus
\[
 \int_G W|u|^2
 \leq\|W\|_p\|u\|_s^2
 \leq \eta\|W\|_p\|\Lcal^{1/2}u\|_2^2
      +C_\eta\|W\|_p\|u\|_2^2.
\]
It follows that the symmetric operator
$T=(\Lcal+V,C_c^\infty(G))$ is bounded from below.

Choose $\lambda>0$ so large that $-\lambda$ lies below this lower bound and
also
\begin{equation}\label{eq:small}
 c_\lambda:=\|W(\Lcal+\lambda)^{-1}\|_{2\to2}<1.
\end{equation}
It is enough, by the semibounded deficiency criterion, to show that
$\ker(T^*+\lambda)=\{0\}$.

Take $f$ in this kernel.  Distributionally,
\[
 \Lcal f+Vf+\lambda f=0.
\]
Since $V\in L^2_{\mathrm{loc}}$ and $f\in L^2$, both $Vf$ and $\Lcal f$
belong to $L^1_{\mathrm{loc}}$.  The sub-Laplacian Kato inequality proved in
\cite[Proposition 1.3]{BC} gives
\begin{equation}\label{eq:kato}
 (\Lcal+\lambda)|f|\leq W|f|
\end{equation}
in the sense of distributions.

Let $0\leq\psi\in C_c^\infty(G)$ and put
$g=(\Lcal+\lambda)^{-1}\psi$.  Positivity of the heat semigroup gives
$g\geq0$; moreover $g\in\Dom(\Lcal)$, $\nabla_{\mathcal H}g\in L^2$, and
hypoellipticity makes $g$ smooth.  To justify testing \eqref{eq:kato} against
$g$, let $\chi_R$ be the standard nonnegative exhausting cutoffs with
\[
 \chi_R\to1,\qquad
 \|\nabla_{\mathcal H}\chi_R\|_\infty\leq C/R,
 \qquad \|\Lcal\chi_R\|_\infty\leq C/R^2.
\]
Test against $\chi_Rg$ and expand
\[
 (\Lcal+\lambda)(\chi_Rg)
 =\chi_R\psi+(\Lcal\chi_R)g
  -2\nabla_{\mathcal H}\chi_R\cdot\nabla_{\mathcal H}g.
\]
The last two terms vanish after pairing with $|f|$ as $R\to\infty$, by
Cauchy--Schwarz and the displayed cutoff bounds.  On the right,
$Wg\in L^2$ by Lemma~\ref{lem:resolvent}, so dominated convergence is also
legitimate.  Hence
\[
 \langle |f|,\psi\rangle
 \leq\langle |f|,W(\Lcal+\lambda)^{-1}\psi\rangle
 \leq c_\lambda\|f\|_2\|\psi\|_2.
\]
Taking the supremum over nonnegative unit vectors in the dense cone
$C_c^\infty(G)_+$ gives
\[
 \|f\|_2\leq c_\lambda\|f\|_2.
\]
By \eqref{eq:small}, $f=0$.  This proves
Theorem~\ref{thm:criterion}.

\section{Application to the source potential}

Let $G$ now be a M\'etivier group, written as
$\R^{2n}\times\R^m$, with homogeneous dimension $Q=2n+2m\geq4$, and let
\[
 N(x,t)=\bigl(|x|^4+16|t|^2\bigr)^{1/4}.
\]
For $w_\alpha=e^{-N^\alpha}$, let $V_\alpha$ be the potential in \cite{BC}.

\begin{corollary}\label{cor:source}
For every $0<\alpha<2$, the operator
\[
 (\Lcal+V_\alpha,C_c^\infty(G))
\]
is essentially self-adjoint.  Consequently the weighted form operator
$\Lcal^{w_\alpha}$ is unitarily equivalent, under
$U_\alpha f=f\sqrt{w_\alpha}$, to the closure of this Schr\"odinger operator.
Thus Remark 4.8 has an affirmative answer in its full stated range.
\end{corollary}

\begin{proof}
Lemma 4.5 of \cite{BC} supplies positive constants $c_{\alpha,1}$ and
$c_{\alpha,2}$ such that, off the identity,
\begin{equation}\label{eq:source-lower}
 V_\alpha(x,t)\geq
 N^{2\alpha-4}|x|^2
 \left(c_{\alpha,1}-\frac{c_{\alpha,2}}{N^\alpha}\right).
\end{equation}
Choose $R$ so that the bracket is nonnegative for $N\geq R$.  Then
$W_\alpha=(V_\alpha)_-$ is supported in $\{N\leq R\}$ and
\begin{equation}\label{eq:negative}
 0\leq W_\alpha(x,t)
 \leq c_{\alpha,2}N^{\alpha-4}|x|^2
 \leq C N^{\alpha-2}.
\end{equation}
Since homogeneous balls have volume $Cr^Q$, \eqref{eq:negative} gives
\[
 W_\alpha\in L^p(G)quad\text{whenever}\quad p(2-\alpha)<Q.
\]
The interval
\[
 \frac Q2<p<\frac{Q}{2-\alpha}
\]
is nonempty exactly because $\alpha>0$.  Furthermore, the exact derivative
bounds in the proof of source Lemma 4.5 give locally
\[
 |V_\alpha|\leq C\bigl(N^{2\alpha-2}+N^{\alpha-2}\bigr).
\]
This belongs to $L^2_{\mathrm{loc}}$ because
$2(2-\alpha)<Q$.  Theorem~\ref{thm:criterion} proves essential
self-adjointness.

It remains only to record the unitary statement, since the weight is not
smooth at the identity.  Put $h=\sqrt{w_\alpha}$.  For
$\phi\in C_c^\infty(G)$,
\[
 h\nabla_{\mathcal H}(h^{-1}\phi)
 =\nabla_{\mathcal H}\phi
  +\frac\alpha2\phi N^{\alpha-1}\nabla_{\mathcal H}N\in L^2.
\]
The last inclusion follows from $Q\geq4$ and $\alpha>0$.  The distributional
product rule, first off the identity and then across it by local
integrability, yields
\[
 \Lcal^{w_\alpha}_1(h^{-1}\phi)
 =h^{-1}(\Lcal+V_\alpha)\phi\in L^2(w_\alpha).
\]
The domain characterization in \cite[Theorem 2.3]{BC} therefore shows that
$h^{-1}\phi\in\Dom(\Lcal^{w_\alpha})$.  Hence
$U_\alpha\Lcal^{w_\alpha}U_\alpha^{-1}$ is a self-adjoint extension of
$(\Lcal+V_\alpha,C_c^\infty)$.  Essential self-adjointness forces it to equal
the closure, completing the proof.
\end{proof}

\section{Verification and novelty}

The exponent-window script checks the strict inequalities, radial
integrability exponent, resolvent decay exponent, and local $L^2$ condition
for the critical dimension $Q=4$ and several higher dimensions.  It is a
sanity check; the proof above is analytic.

A bounded search through 13 August 2026 covered the run indexes, the exact
open phrase, arXiv:1610.09850, the paper title, M\'etivier Schr\"odinger
potentials, and essential self-adjointness.  It found the source and later
papers citing it, but no explicit answer to Remark 4.8.  This is not an
exhaustive priority search.

\paragraph{Human-review recommendation.}
Check first the cutoff passage from \eqref{eq:kato} to the resolvent test, and
then the distributional conjugation across the identity in the final
paragraph.  All exponent inequalities are strict for every $\alpha>0$.

\begin{thebibliography}{9}
\bibitem{BC}
T. Bruno and M. Calzi,
\emph{Weighted sub-Laplacians on M\'etivier groups: essential
self-adjointness and spectrum},
Proc. Amer. Math. Soc. 145 (2017), 3579--3594;
arXiv:1610.09850.

\bibitem{Folland}
G. B. Folland,
\emph{Subelliptic estimates and function spaces on nilpotent Lie groups},
Ark. Mat. 13 (1975), 161--207.

\bibitem{RS}
M. Reed and B. Simon,
\emph{Methods of Modern Mathematical Physics II: Fourier Analysis,
Self-Adjointness}, Academic Press, 1975.
\end{thebibliography}

\end{document}
